\section{Introduction}
\label{sec:introduction}

The multilayer perceptron (MLP) is where deep learning starts.
Stack a few affine maps, put a nonlinearity between them, and train the whole
thing by backpropagation~\citep{rumelhart1986}: that recipe is still the first
model most people reach for, and it underlies the feedforward blocks inside
larger architectures~\citep{goodfellow2016}.

Theory has long said the MLP is expressive.
A single hidden layer, made wide enough, approximates any continuous function on
a compact set~\citep{cybenko1989, hornik1989}.
What the theory does not say is \emph{how} wide, or whether the width that fits
the training data is the width that generalizes.
That gap between existence and practice is the question we poke at here.

This example paper sets up an MLP (\cref{sec:model}, sketched in
\cref{fig:arch}), trains it with the procedure in \cref{alg:train}
(\cref{sec:training}), and studies the effect of hidden width in
\cref{sec:experiments}.
The finding is unsurprising but worth seeing plotted: accuracy climbs steeply
with width and then flattens, so the useful capacity is far smaller than the
approximation theorems would let you use.
